
\textbf{Weight-space learning.} Hyper-representations are learned embeddings of neural network weights used for transfer learning tasks, demonstrating that weight-based features generalize across model zoos. Model stitching analyses show that layers can be examined semi-independently while preserving model behavior, supporting the hypothesis that cross-layer dependencies are not always critical. Weight statistics such as norms and spectral properties have been shown to correlate with model accuracy. Building on this prior validation of weight-space signal, the present work compares statistical aggregation directly against transformer-based tokenization, finding that on small datasets (100 models), both approaches achieve 80\% accuracy.

\textbf{Neural functionals.} Recent methods process weights as graph structures or functional mappings. While these approaches show promise, they introduce architectural complexity. The present work isolates tokenization quality independent of backbone architecture choice, providing a baseline for future comparisons.

\textbf{Sequence modeling on structured data.} Point cloud transformers and graph transformers process variable-length structured data via attention mechanisms. Weight-space learning shares similar challenges---variable layer dimensions and permutation symmetries---but operates on learned parameter distributions rather than geometric coordinates. Hierarchical position encodings capture depth information in trees and nested structures. This work adapts such encodings for weight tokens where layer depth indicates network position.

\textbf{Meta-learning on model zoos.} Neural architecture search methods predict model performance from architectural descriptors. Weight-space learning extends this by using learned parameters as features rather than discrete architecture choices, enabling fine-grained property prediction post-training. Model merging methods combine pretrained models via weight averaging or interpolation. Tokenization is prerequisite for learning how to merge---embeddings must preserve model properties to predict merge outcomes. Backdoor detection methods use spectral signatures and statistical anomalies in weight distributions. The present validation that layer-wise tokenization retains sufficient signal (80\% family classification) suggests discriminative capacity for finer-grained backdoor patterns.

\textbf{Positioning.} Existing weight-space learning methods demonstrate task-specific success but lack systematic evaluation of tokenization strategies. The present work provides a controlled comparison of layer-wise processing approaches, showing that (1) tokenization preserves signal (80\% accuracy exceeds random baseline 25\% and threshold 60\%), (2) simple baselines matter (per-layer statistics match transformer performance on small datasets), and (3) family separability dominates (architecture-specific patterns are detectable via local layer statistics without cross-layer reasoning). This validation enables future work on larger datasets and more complex tasks with confidence that the layer-wise assumption holds.
